\documentclass[10pt,a4paper]{amsart}
\usepackage[margin=19mm]{geometry}
\usepackage{amsmath,amssymb,mathtools}
\usepackage[scr=rsfs,cal=euler]{mathalfa}
\usepackage{xfrac}
\usepackage[unicode,hidelinks,bookmarks=false]{hyperref}
\newcommand{\ie}{i.e.\ }
\newcommand{\cf}{cf.\ }
\newcommand{\resp}{respectively\ }
\newcommand{\Subsection}{\subsection}
\providecommand{\nobreakdash}{\nobreak\hbox{-}}
\setlength{\emergencystretch}{2em}
\multlinegap0pt
\usepackage{bm}
\usepackage{bbm}
\usepackage{dsfont}
\usepackage{xspace}
\usepackage{cancel}
\usepackage{mathtools}
\usepackage{amsthm}
\makeatletter
\@ifundefined{theorem}{\newtheorem{theorem}{Theorem}}{}
\@ifundefined{proposition}{\newtheorem{proposition}[theorem]{Proposition}}{}
\makeatother
\def\E{\mathbb E}
\def\P{\mathbb P}
\def\TT{\widetilde{T}}
\title[Distribution of random multiplicative functions in short int]{Distribution of random multiplicative functions in short intervals, with proper normalization}
\author{Adam J. Harper, Kannan Soundararajan, Max Wenqiang Xu}
\date{Source: arXiv:2606.29040v1 (CC BY 4.0)}
\begin{document}
\maketitle

\noindent\emph{Source and licence.} Distribution of random multiplicative functions in short intervals, with proper normalization. Version arXiv:2606.29040v1 (CC BY 4.0), distributed under the Creative Commons Attribution 4.0 International licence, as recorded in the arXiv licence metadata for this exact version and shown on the abstract page https://arxiv.org/abs/2606.29040v1. Full rights notice: licenses/arxiv-2606.29040-CC-BY-4.0.txt.\par\smallskip

\noindent\textit{Editorial scope.} Counted unit: the Proposition whose source label is \texttt{prop3.3} in the article (source0.tex lines 648--652, source label \texttt{prop3.3}), delivered together with the article's own complete original proof of it (source0.tex lines 653--718, one \texttt{proof} environment of 66 lines). No mathematics is written, repaired or re-derived: statement and proof are the article's own text line for line. Required context retained uncounted: 3.4 (lines 624--627); eqn: g(t) (lines 630--632). Every definition, display and label of those lines is retained unchanged. Omitted from this package and not redistributed: 1--623, 628--629, 633--647, 719--1780. Nothing inside a retained excerpt is omitted or altered: every retained body is delivered exactly as the article's lines are written, so no omission receipt of any kind accompanies this package. Citations: the retained excerpts carry no citation command at all, so no bibliography entry of the article is transcribed and no reference is redistributed. Reference adaptation: none was needed and none was made; every reference inside the delivered unit resolves inside it, so no unresolved reference or citation is produced. Printed slips kept literal: no printed word, symbol, hypothesis or number of the counted statement or of its proof was altered. Layout and class: the article's own class is \texttt{amsart} with packages including latexsym, amsfonts, amsmath, amssymb, amsthm, cite, graphicx, hyperref, bbm, amscd, tikz, mathrsfs, url, inputenc; the wrapper is the lane's pinned \texttt{amsart} editorial wrapper at 10pt on A4, which restates the theorem-like environments the retained text needs and adds the article's own macro definitions that the retained text needs (\texttt{\textbackslash E}, \texttt{\textbackslash P}, \texttt{\textbackslash TT}), with extra wrapper packages (bm, bbm, dsfont, xspace, cancel, mathtools). The delivered numbering is the unit's own. Assets: no retained span contains a figure environment, a graphics-inclusion command, a PDF-page inclusion, a table or a TikZ picture; no image file is copied into this package, the package declares no asset dependency and no image file is redistributed with it. Licence: version arXiv:2606.29040v1 of this article is distributed under the Creative Commons Attribution 4.0 International licence, as recorded in the arXiv licence metadata for this exact version and shown on the abstract page https://arxiv.org/abs/2606.29040v1. A full rights notice is delivered at licenses/arxiv-2606.29040-CC-BY-4.0.txt. Characters: every retained excerpt is pure ASCII, so the delivered engine, XeLaTeX with its default Latin Modern text font, reports no missing character.
 \begin{equation} 
 \label{3.4} 
 \E[ |{\mathcal F}_j(1/2+it)|^2 ] = \prod_{z^{e^{-\tau}} < p \le z^{e^{-j}}} \Big( 1- \frac{1}{p}\Big)^{-1} \sim e^{\tau-j}.
 \end{equation} 

\begin{equation}\label{eqn: g(t)}
     (\sqrt{T} e^{\tau-j})^{-B}  \le |{\mathcal F}_j(\tfrac 1 2 +it)|  \le \sqrt{T} e^{\tau-j}\exp((\log T)^{\frac 1{100}}) (\log \log z - j).   
\end{equation}

\begin{proposition} \label{prop3.3}  Recalling that $\TT = T(\log T)^{100}$, we have  
\begin{equation}\label{eqn: exceptprob}
\P\Big( \mathcal{G}(t)~\text{holds for all $|t|\le \TT$} \Big)  \ge 1 -\exp (-(\log T)^{\frac 1{100}}). 
\end{equation}
\end{proposition} 

\begin{proof}   Given $0\le j\le \tau -1$  we shall show that the probability that \eqref{eqn: g(t)} fails for some $|t| \le \TT$ is  
$$ 
\ll \exp(-(\log T)^{\frac 1{100}}) (\log \log z -j)^{-2}.  
$$
Summing this over all the possibilities for $j$ yields the proposition.  

Consider a mesh of points ${\mathcal T}_j = \{ {\widehat t} = e^{j} n/\log z :\ \ n\in {\mathbb Z}, \ |{\widehat t}| \le \TT\}$.  The mesh ${\mathcal T}_j$ contains $\ll \TT e^{-j} \log z$ points, and for each $t$ with $|t|\le \TT$ we may find ${\widehat t} \in {\mathcal T}_j$ with $|t-{\widehat t}| \le e^{j}/\log z$. It will turn out that this places ${\widehat t}$ sufficiently close to $t$ (relative to the length of the Euler product ${\mathcal F}_{j}$) that the behaviour of ${\mathcal F}_j(\tfrac 12 +it)$ is essentially controlled by that of ${\mathcal F}_j(\tfrac 12 +i{\widehat t})$.

Indeed, if \eqref{eqn: g(t)} fails at $t$, then with ${\widehat t}$ denoting the nearest point to $t$ in ${\mathcal T}_j$ we must have one of the following four possibilities: 
\begin{equation} 
\label{3.8} 
|{\mathcal F}_j(\tfrac 12 +i{\widehat t})| \ge \tfrac 12 \sqrt{T} e^{\tau-j} \exp((\log T)^{\frac 1{100}}) (\log \log z - j), 
\end{equation} 
or 
\begin{equation} 
\label{3.9} |{\mathcal F}_j(\tfrac 12 + i{\widehat t})|^{-1} \ge \tfrac 12 (\sqrt{T} e^{\tau-j})^B, 
\end{equation} 
or 
\begin{equation} 
\label{3.10} 
\int_{-e^{j}/\log z}^{e^{j}/\log z} |{\mathcal F}_j^{\prime}(\tfrac 12+i{\widehat t} +ih)| dh 
\ge |{\mathcal F}_j(\tfrac 12+it)-{\mathcal F}_j(\tfrac 12+i{\widehat t})| \ge \tfrac 12 \sqrt{T} e^{\tau-j} \exp((\log T)^{\frac 1{100}}) (\log \log z - j), 
\end{equation}
or 
\begin{equation} 
\label{3.11} 
\int_{-e^{j}/\log z}^{e^{j}/\log z} |({\mathcal F}_j(\tfrac 12+i{\widehat t} +ih)^{-1})^{\prime}| dh 
\ge  |{\mathcal F}_j(\tfrac 12+it)^{-1} - {\mathcal F}_j(\tfrac 12 +i{\widehat t})^{-1} |  > \tfrac 12  (\sqrt{T}e^{\tau-j})^{B}. 
\end{equation} 

Given ${\widehat t}$, using Markov's inequality with \eqref{3.4}, we see that the probability that \eqref{3.8} holds is
\begin{align} 
\label{3.12} 
&\ll e^{\tau-j} (\tfrac 12 \sqrt{T} e^{\tau-j} \exp((\log T)^{\frac 1{100}}) (\log \log z - j))^{-2} \nonumber\\
& \ll T^{-1} e^{j-\tau} \exp(-2(\log T)^{\frac 1{100}}) (\log \log z - j)^{-2}.
\end{align} 
Since 
$$
\E[|{\mathcal F}_j(\tfrac 12+i {\widehat t})|^{-2}] = \prod_{z^{e^{-\tau}} < p \le z^{e^{-j}}} \Big(1 +\frac 1p \Big) \sim e^{\tau -j}, 
$$ 
an even stronger bound applies for the probability that \eqref{3.9} holds.   Next, note that 
$$ 
\E[ |{\mathcal F}_j^{\prime}(\tfrac 12 +i{\widehat t} +ih)|^2 ] = \E\Big[ \Big| \sum_{\substack{ n=1 \\ p|n \implies z^{e^{-\tau}}<p\le z^{e^{-j}}}} \frac{f(n) \log n}{n^{\frac 12+i {\widehat t}+ih}}\Big|^2\Big] = \sum_{\substack{ n=1 \\ p|n \implies z^{e^{-\tau}}<p\le z^{e^{-j}}}} \frac{(\log n)^2}{n},  
$$ 
and using $(\log n)/n^{\alpha} \le 1/(e\alpha)$ for all $\alpha>0$ and $n\ge 1$ we may bound this by 
\begin{align*} 
&\ll \Big( \frac{\log z}{e^j} \Big)^2 \sum_{\substack{ n=1 \\ p|n \implies z^{e^{-\tau}}<p\le z^{e^{-j}}}} \frac{1}{n^{1-2 e^{j}/\log z}} 
\\
&\le \Big( \frac{\log z}{e^j} \Big)^2  \prod_{z^{e^{-\tau}}<p\le z^{e^{-j}}} \Big(1- \frac 1{p^{1-2e^{j}/\log z}}\Big)^{-1} \ll \Big( \frac{\log z}{e^j} \Big)^2 e^{\tau-j}. 
\end{align*} 
Using these estimates and Cauchy--Schwarz we conclude that 
$$ 
\E \Big[ \Big| \int_{-e^{j}/\log z}^{e^{j}/\log z} |{\mathcal F}_j^{\prime}(\tfrac 12+i{\widehat t} +ih)| dh \Big|^2 \Big] 
\ll \frac{e^{j}}{\log z} \int_{-e^{j}/\log z}^{e^{j}/\log z}  \E[ |{\mathcal F}_j^{\prime}(\tfrac 12 +i{\widehat t} +ih)|^2 ]  dh \ll e^{\tau- j}. 
$$ 
Therefore by Markov's inequality, the probability that \eqref{3.10} holds is also bounded by the quantity in \eqref{3.12}.   An entirely analogous argument shows that the same estimate also holds for the probability with which \eqref{3.11} holds.   

Since there are $\ll {\widetilde T}e^{-j} \log z$ possible points ${\widehat t}$ in ${\mathcal T}_j$, we conclude that the probability that one of the four possibilities in \eqref{3.8}, 
\eqref{3.9},  \eqref{3.10}, or \eqref{3.11} holds for some ${\widehat t} \in {\mathcal T}_j$ is 
\begin{align*}
&\ll ({\widetilde T} e^{-j} \log z) T^{-1} e^{j-\tau} \exp(-2(\log T)^{\frac 1{100}}) (\log \log z - j)^{-2}\\
& \ll (\log T)^{100} (e^{-\tau} \log z) \exp(-2(\log T)^{\frac 1{100}}) (\log \log z - j)^{-2} \\
&\ll \exp(-(\log T)^{\frac 1{100}}) (\log \log z - j)^{-2}, 
\end{align*}
upon recalling the definition of $\tau$.   This bounds the probability that \eqref{eqn: g(t)} fails for this particular $j$ and some $|t| \le \TT$, and completes our proof. 
 \end{proof} 

\end{document}
